\documentclass[11pt,a4paper]{article}
\usepackage{ctex}
\usepackage{amsmath,amssymb}
\usepackage{booktabs}
\usepackage{graphicx}
\usepackage[margin=2.5cm]{geometry}
\usepackage{algorithm}
\usepackage{algorithmic}
\usepackage{hyperref}

\title{是语切词：一种完全非监督的切词器构建方法}
\author{}
\date{}

\begin{document}
\maketitle

\begin{abstract}
中文分词的质量在很大程度上取决于词表——什么算词、包含哪些词。
本文提出一种完整可复现的分词方案：以中文维基百科、维基词典等公开知识库作为词表来源，以维基百科语料作为训练数据，通过EM迭代学习词频，最终得到一个Unigram分词模型。
整个流程从数据获取到词表训练全部自动化，不需要任何人工标注。
我们认为，维基体系作为词表来源具有三个关键优势：标准统一、覆盖面广、持续更新。
系统以不到1000行C++实现，完整代码开源。
\end{abstract}

\section{引言}

中文没有显式的词边界，分词是几乎所有中文NLP任务的基础步骤。
围绕分词问题，学术界已经发展出了从规则方法到深度学习的一系列技术。
然而，一个常被忽视的核心问题是：\textbf{分词的好坏，除了边界判断的准确性，很大程度上取决于分词标准本身——什么是词，词表里应该包含哪些词。}

不同的标注规范（如北大标准、微软标准）对"词"的定义不同，导致同一句话的"正确"分词结果也不同。
监督方法通过拟合特定标注数据来回避这个问题，但标注规范的选择本身就是一个人为决策，且标注数据昂贵、领域受限。

我们注意到，维基百科和维基词典构成了一个天然的、高质量的词表来源：

\begin{itemize}
    \item \textbf{标准统一。} 维基百科的收录标准经过社区长期讨论和共识，对"什么值得成为一个词条"有清晰的判断。维基词典则专注于语言学意义上的词汇收录。这为"什么是词"提供了一个客观、一致的标准。
    \item \textbf{覆盖面广。} 中文维基百科包含超过130万个词条，维基词典收录了大量日常词汇和成语。两者互补，从专有名词到通用词汇都有覆盖。
    \item \textbf{持续更新。} 与静态的标注数据集不同，维基百科每天都在更新，新词汇、新实体会被不断收录。这意味着基于维基词表的分词器可以通过简单地重新获取词表来跟上语言的演化，而不需要重新标注训练数据。
\end{itemize}

基于这一观察，本文提出Iscut系统：从维基体系获取词表，从维基百科语料学习词频，通过EM迭代训练一个Unigram分词模型。
整个流程完全自动化、可复现，不依赖任何人工标注。

\section{系统概览}

Iscut的完整流程分为四个阶段：

\begin{enumerate}
    \item \textbf{词表获取}（\texttt{conv/}）：从Wikimedia下载中文维基百科、维基词典、维基文库的标题数据和重定向数据，经繁简转换和规则过滤，生成候选词表。
    \item \textbf{语料获取}（\texttt{data/}）：从HuggingFace下载中文维基百科和维基新闻的文本数据，经繁简转换后合并为训练语料。
    \item \textbf{词频训练}（\texttt{scripts/}）：统计语料字频以过滤含稀有字的词，然后通过EM迭代学习词频分布。
    \item \textbf{分词}（\texttt{src/}）：基于学到的词频，用DAG+动态规划进行分词。
\end{enumerate}

每个阶段都通过Makefile自动化，支持单步执行或一键运行。

\section{词表获取}

\subsection{数据来源}

我们从以下四个来源获取候选词：

\begin{itemize}
    \item \textbf{中文维基百科}（zhwiki）：所有正文页面的标题。
    \item \textbf{中文维基词典}（zhwiktionary）：所有词条标题，主要包含日常词汇和成语。
    \item \textbf{中文维基文库}（zhwikisource）：所有文献标题。
    \item \textbf{维基百科重定向}：从SQL dump中提取的重定向映射，重定向的源标题和目标标题都作为候选词。重定向数据特别有价值，因为它包含了大量同义词和别名（如"电脑"重定向到"计算机"）。
    \item \textbf{网络用语}：从维基百科"中国大陆网络用语列表"页面提取的流行词汇。
\end{itemize}

所有来源均通过OpenCC进行繁简转换。

\subsection{过滤规则}

候选词经过以下规则过滤：

\begin{itemize}
    \item 所有词必须为纯汉字。
    \item 来自维基词典的词保留$\leq 4$字的（4字词大概率为成语）。
    \item 其它来源保留$\leq 3$字的。
    \item 维基百科中含间隔号"·"的词条（如"安妮·莎士比亚"）按"·"拆分，每段$\leq 4$字的保留。这有效地处理了大量外国人名和地名。
    \item 过滤掉包含"消歧义""列表"等子串的词条。
\end{itemize}

此外，在训练阶段还会进一步用语料字频过滤：词中每个字都必须在语料中出现至少2次，以去除包含极稀有字的词。

\section{语料获取与处理}

训练语料来自两个HuggingFace数据集：

\begin{itemize}
    \item \textbf{FineWiki zhwiki}（HuggingFaceFW/finewiki的中文子集）：约130万篇维基百科文章，经过清洗的高质量文本。
    \item \textbf{roots\_zh\_wikinews}（bigscience-data/roots\_zh\_wikinews）：中文维基新闻，提供新闻领域的补充。
\end{itemize}

处理步骤：
\begin{enumerate}
    \item 从JSONL中提取文本字段。
    \item 去除Markdown标题标记。
    \item OpenCC繁简转换。
    \item 每篇文章合并为一行，空白折叠。
    \item 两个来源合并为单一语料文件。
\end{enumerate}

\section{分词方法}

\subsection{标点预切分}

在进行DAG+DP分词之前，先按标点符号将输入切分为若干子句。
标点符号本身作为独立的token输出，每个子句独立进行分词。
这一步有两个好处：一是避免标点参与跨标点的词匹配，二是缩短每段的长度以提高效率。

标点的判断覆盖ASCII标点、CJK标点符号区（U+3000--303F）、全角标点（U+FF01--FF65）和通用标点区（U+2000--206F）。

\subsection{冷启动：正向最长匹配}

在没有词频信息时，用正向最长匹配（Forward Maximum Matching）初始化。
对输入的每个位置，在词典中查找最长匹配词；若无匹配，则输出单个字符并前进一位。
FMM倾向于选择长词，为第一轮EM提供初始词频估计。

\subsection{基于双数组Trie的DAG构建}

词典存储在双数组Trie中，支持高效的公共前缀搜索。
对句子$s$，构建有向无环图（DAG）：
\begin{itemize}
    \item 节点对应字节位置$0, 1, \ldots, n$。
    \item 对每个位置$i$，执行前缀搜索找到所有匹配词$w \in \mathcal{D}$，每个匹配产生一条边，权重为$\log P(w)$。
    \item 始终保留单字符兜底边。
\end{itemize}

\subsection{动态规划求最优路径}

在DAG上用后向动态规划求最大权路径：
\begin{equation}
    V(i) = \max_{(i,j) \in E} \left[ \log P(w_{ij}) + V(j) \right]
\end{equation}
边界条件$V(n) = 0$。通过回溯即可恢复最优分词。

\subsection{EM迭代}

EM循环交替执行两步：

\textbf{E步}（分词）：给定当前词频$\{f(w)\}$，计算$P(w) = f(w) / \sum_{w'} f(w')$，用DAG+DP对全部语料分词。

\textbf{M步}（计数）：从分词结果重新统计词频：
\begin{equation}
    f^{(t+1)}(w) = \sum_{s \in \mathcal{C}} \text{count}(w, \text{seg}^{(t)}(s))
\end{equation}

只统计出现在原始词典$\mathcal{D}$中的词，保持词表在迭代间稳定。

\begin{algorithm}[t]
\caption{基于固定词典的EM分词算法}
\begin{algorithmic}[1]
\REQUIRE 词典$\mathcal{D}$，语料$\mathcal{C}$，迭代轮数$T$
\STATE 用$\mathcal{D}$构建Trie
\STATE $\text{seg} \leftarrow \text{FMM}(\mathcal{C}, \mathcal{D})$ \COMMENT{冷启动}
\STATE $f \leftarrow \text{Count}(\text{seg}, \mathcal{D})$ \COMMENT{初始词频}
\FOR{$t = 1$ to $T$}
    \STATE $P(w) \leftarrow f(w) / \sum_{w'} f(w')$
    \STATE $\text{seg} \leftarrow \text{DAG-DP}(\mathcal{C}, P)$ \COMMENT{E步}
    \STATE $f \leftarrow \text{Count}(\text{seg}, \mathcal{D})$ \COMMENT{M步}
\ENDFOR
\RETURN $f$
\end{algorithmic}
\end{algorithm}

\subsection{收敛性}

定义联合目标函数：
\begin{equation}
    F(\theta, Z) = \sum_{m=1}^{M} \sum_{w \in z_m} \log P_\theta(w)
\end{equation}

每轮迭代中，E步对$Z$做$\arg\max$，M步对$\theta$做最大似然估计，两步都不降低$F$的值。
由于词典和语料有限，可能的分词方案构成有限集，因此$F$的单调不减序列必在有限步内收敛。

\section{实现}

系统以C++17实现，无外部依赖，核心代码不到1000行。

\begin{itemize}
    \item 双数组Trie采用XOR索引，公共前缀搜索复杂度为$O(L)$。
    \item 全部文本处理在UTF-8字节序列上进行，避免码点转换开销。
    \item 词频以整数存储，分词时转换为对数概率。
    \item 提供Python绑定（pybind11），可通过\texttt{pip install}安装使用。
\end{itemize}

整个训练流程通过Makefile自动化：

\begin{verbatim}
cd conv && make          # 获取并过滤维基词表
cd data && make          # 获取并处理语料
cd scripts && make       # 字频过滤 + EM训练 + 替换词典
\end{verbatim}

\section{讨论}

\subsection{为什么用维基词表}

传统的分词方法要么依赖人工编纂的词典（规模有限、更新缓慢），要么依赖标注数据隐式定义词表（昂贵、领域受限）。
维基体系提供了一个独特的替代方案：

\begin{enumerate}
    \item \textbf{词表即知识。}
    分词标准的核心是"什么是词"的判断。维基百科的收录标准经过全球社区的长期讨论，维基词典则由语言学爱好者维护。
    两者共同提供了一个经过广泛共识的词表标准。

    \item \textbf{活的词表。}
    语言是不断演化的。新的人名、地名、术语、流行语每天都在产生。
    维基百科作为世界上更新最活跃的知识库之一，能够及时收录这些新词。
    基于维基词表的分词器只需重新运行获取流程，就能自动包含最新的词汇，而不需要重新标注任何训练数据。

    \item \textbf{多源互补。}
    维基百科覆盖专有名词和实体，维基词典覆盖日常词汇和成语，维基文库提供古典文献中的词汇，重定向数据补充同义词和别名。
    这种多源组合提供了比任何单一词典更全面的覆盖。
\end{enumerate}

\subsection{为什么不用软EM}

原则上可以用前向-后向算法计算所有可能分词上的期望计数（软EM）。
我们选择硬EM是出于简洁性和效率的考虑。
从实验看，硬EM收敛良好，MAP分词已捕获了大部分概率质量。

\subsection{局限性}

\begin{itemize}
    \item 长度过滤规则是启发式的，可能丢弃部分有价值的长词。
    \item 当前不处理未登录词的发现，只能对未登录字做单字符兜底。
    \item Unigram模型不考虑上下文，对歧义的处理能力有限。
\end{itemize}

\section{结论}

本文提出Iscut，一个完整可复现的中文分词系统。
其核心观点是：维基百科和维基词典构成了一个天然的、持续更新的高质量词表来源，将其与EM词频学习相结合，可以在不需要任何人工标注的情况下得到一个实用的分词器。
整个流程从数据获取到训练完全自动化，代码开源，任何人都可以一键复现。


\end{document}
